\chapter{Formulation}
\label{ch:formulation}
\impl{src/physics.jl}

\section{Strong and weak form}

Carina solves finite-deformation solid mechanics in the reference
configuration.  Linear momentum balance reads
\begin{equation}
  \Grad \cdot \Piola + \rho_0 \vect{b} = \rho_0 \ddot{\disp},
  \label{eq:momentum}
\end{equation}
with $\Piola$ the first Piola--Kirchhoff stress, $\vect{b}$ a body force per
unit mass, and $\rho_0$ the reference density.  Multiplying by an admissible
variation $\delta\disp$ vanishing on the Dirichlet boundary $\Gamma_u$ and
integrating by parts gives the weak form: find $\disp$ with $\disp =
\bar{\disp}$ on $\Gamma_u$ such that
\begin{equation}
  \int_{\Omega_0} \Piola \dbldot \Grad \delta\disp \, \mathrm{d}V
  + \int_{\Omega_0} \rho_0\, \ddot{\disp} \cdot \delta\disp \, \mathrm{d}V
  = \int_{\Omega_0} \rho_0 \vect{b}\cdot\delta\disp \, \mathrm{d}V
  + \int_{\Gamma_t} \bar{\vect{t}} \cdot \delta\disp \, \mathrm{d}A
  \label{eq:weak}
\end{equation}
for all admissible $\delta\disp$, where $\bar{\vect{t}}$ is the prescribed
traction on $\Gamma_t$.

The kinematics enter only through the deformation gradient
\begin{equation}
  \Fdef = \Ident + \Grad\disp ,
\end{equation}
and the constitutive response enters only through $\Piola(\Fdef)$.  This is the
entire coupling surface between Chapter~\ref{ch:constitutive} and everything
downstream: a constitutive model supplies $\Piola$ given $\Grad\disp$, and
optionally its derivative.

\section{Discretization}

Introducing isoparametric shape functions $N_a$ over an element mesh and
writing $\disp \approx \sum_a N_a \Uvec_a$, \eqref{eq:weak} becomes the
semi-discrete system
\begin{equation}
  \Mmat \Avec + \Fint(\Uvec) = \Fext(t),
  \label{eq:semidiscrete}
\end{equation}
with
\begin{equation}
  \Mmat_{ab} = \int_{\Omega_0} \rho_0 N_a N_b \, \mathrm{d}V,
  \qquad
  \Fint_a(\Uvec) = \int_{\Omega_0} \Piola(\Fdef) \cdot \Grad N_a \, \mathrm{d}V .
\end{equation}
Both integrals are evaluated by Gauss quadrature element by element.  The
residual that every solver in this manual drives to zero is
\begin{equation}
  \Rvec(\Uvec) = \Fext - \Fint(\Uvec) - \Mmat\Avec ,
  \label{eq:residual}
\end{equation}
with the inertial term absent in the quasi-static case
(Section~\ref{sec:quasistatic}).

\begin{remark}
Carina stores $\Rvec$ already negated relative to the force residual, so the
Newton system is $\Kmat\,\Delta\Uvec = \Rvec$ with a plus sign throughout the
implementation.  The sign convention is stated in \code{linear\_solvers.jl} at
the dispatch point and is followed uniformly here.
\end{remark}

\section{Consistent tangent}

Differentiating \eqref{eq:residual} with respect to the nodal displacements
gives the tangent stiffness
\begin{equation}
  \Kmat_{ab} = \frac{\partial \Fint_a}{\partial \Uvec_b}
             = \int_{\Omega_0} \Grad N_a \cdot \fourth{A} \cdot \Grad N_b \,
               \mathrm{d}V,
  \qquad
  \fourth{A} = \frac{\partial \Piola}{\partial \Grad\disp} .
  \label{eq:tangent}
\end{equation}
The fourth-order material tangent $\fourth{A}$ is the only object a
constitutive model must supply beyond the stress itself; how Carina obtains it
--- in closed form, by automatic differentiation, or by finite differences ---
is the subject of Section~\ref{sec:tangent-strategies}.

\section{Dirichlet conditions by elimination}
\label{sec:elimination}

Prescribed displacements are imposed by removing the corresponding equations
and unknowns, not by adding a penalty term or a Lagrange multiplier.
Partitioning the degrees of freedom into unknowns $\mathrm{f}$ and prescribed
values $\mathrm{p}$,
\begin{equation}
  \begin{bmatrix}
    \Kmat_{\mathrm{ff}} & \Kmat_{\mathrm{fp}} \\
    \Kmat_{\mathrm{pf}} & \Kmat_{\mathrm{pp}}
  \end{bmatrix}
  \begin{bmatrix} \Delta\Uvec_{\mathrm{f}} \\ \Delta\Uvec_{\mathrm{p}} \end{bmatrix}
  =
  \begin{bmatrix} \Rvec_{\mathrm{f}} \\ \Rvec_{\mathrm{p}} \end{bmatrix},
\end{equation}
only the first block row is solved, with $\Delta\Uvec_{\mathrm{p}}$ known.  The
cross term $\Kmat_{\mathrm{fp}}\Delta\Uvec_{\mathrm{p}}$ is carried
automatically because the element kernels read displacement fields that already
hold the prescribed values at constrained nodes; nothing assembles
$\Kmat_{\mathrm{fp}}$ explicitly.

This choice is not cosmetic.  A penalty imposition multiplies selected
diagonal entries by a large factor --- $10^{15}$ is typical --- which raises
the condition number of the operator by the same factor and destroys the
convergence rate of every Krylov method in Chapter~\ref{ch:krylov}.  The
eliminated system carries no artificial conditioning at all, so the iteration
counts reported throughout this manual reflect the physics of the problem
rather than the bookkeeping of the boundary conditions.

\begin{remark}
One place in Carina does see penalty-scale entries: the assembled Chebyshev
path, which forms $\Kmat$ including constrained rows before scaling.  That is
why the polynomial there is built on the symmetrically scaled operator
$\Dmat^{-1/2}\Kmat\Dmat^{-1/2}$, whose eigenvalues are $O(1)$ even when
$\Kmat$'s are not.  See Section~\ref{sec:chebyshev}.
\end{remark}

\section{Mass matrix and lumping}

The consistent mass of \eqref{eq:semidiscrete} is used by the implicit dynamic
integrator.  The explicit integrator instead uses the row-sum lumped mass
\begin{equation}
  \Mmat^{\mathrm{L}}_{aa} = \sum_b \Mmat_{ab},
\end{equation}
which is diagonal and therefore invertible in $O(n)$ with no linear solve at
all --- the property that makes explicit integration attractive
(Section~\ref{sec:central-difference}).
